\section{Impactos sociales y ambientales del proyecto}

Más allá de su viabilidad económica, ParkLink tiene el potencial de generar beneficios concretos sobre el entorno urbano, la calidad de vida de los ciudadanos y el medioambiente. Estos efectos, aunque difíciles de monetizar directamente, constituyen una parte fundamental de la propuesta de valor del proyecto.

\subsection*{Reducción de la congestión vehicular}

Una fracción significativa del tránsito urbano en ciudades densas no obedece a desplazamientos con destino definido, sino a vehículos circulando en búsqueda de estacionamiento. Estudios realizados en distintas ciudades del mundo estiman que entre el 30\% y el 40\% del tráfico en zonas céntricas corresponde a conductores buscando dónde estacionar \autocite{shoup2006, shoup2005}. En CABA, con aproximadamente 1,8 millones de vehículos circulando diariamente, incluso una mejora marginal en la eficiencia de este proceso tendría un impacto agregado considerable sobre la fluidez del tránsito.

Al permitir que el conductor reserve su lugar antes de salir, ParkLink elimina la circulación improductiva: el vehículo se desplaza directamente hacia el espacio confirmado, sin vueltas ni tiempos muertos. A escala, esto implica menos autos en movimiento sin destino claro, menor ocupación de arterias viales y una mejora tangible en los tiempos de traslado.

\subsection*{Disminución de emisiones contaminantes}

La circulación innecesaria en busca de estacionamiento no solo genera congestión; también implica consumo adicional de combustible y emisión de gases de efecto invernadero. Cada minuto de motor encendido sin avance productivo equivale a combustible quemado y dióxido de carbono liberado sin contrapartida de utilidad.

Si ParkLink logra evitar, en promedio, entre 10 y 15 minutos de búsqueda por reserva exitosa, el ahorro acumulado en emisiones a lo largo de miles de operaciones diarias resultaría cuantificable y significativo. En un contexto global donde las ciudades buscan reducir su huella de carbono y mejorar la calidad del aire, la optimización del estacionamiento urbano representa una palanca concreta de intervención.

\subsection*{Mejor aprovechamiento de la infraestructura existente}

Uno de los problemas identificados en el análisis es la subutilización de espacios privados que permanecen vacíos durante buena parte del día. ParkLink no propone construir nueva infraestructura, sino hacer visible y accesible la que ya existe. Esto se alinea con el principio de economía colaborativa: extraer valor de activos ociosos, reduciendo la presión sobre los recursos urbanos disponibles y postergando la necesidad de ampliar la oferta de estacionamiento público.

\subsection*{Beneficios sobre la calidad de vida}

La búsqueda de estacionamiento es una fuente frecuente de estrés, demoras y frustración para los conductores urbanos. La posibilidad de planificar el estacionamiento con anticipación reduce la incertidumbre del traslado, mejora la experiencia del conductor y libera tiempo que de otro modo se destinaría a circular sin rumbo. Este beneficio, aunque subjetivo, tiene impacto directo sobre el bienestar cotidiano de una porción relevante de la población activa.

Adicionalmente, para los propietarios de cocheras, ParkLink representa la posibilidad de generar un ingreso pasivo a partir de un bien inmueble subutilizado, mejorando su situación económica sin requerir esfuerzo operativo adicional.

\subsection*{Alineación con políticas de movilidad urbana sostenible}

Las tendencias actuales en planificación urbana priorizan la reducción del uso del automóvil privado y la mejora de la eficiencia del tránsito. Si bien ParkLink no desincentiva el uso del auto, sí contribuye a que cada viaje sea más eficiente, reduciendo la circulación innecesaria que deriva de la falta de información sobre disponibilidad. En este sentido, la plataforma es compatible con y complementaria a políticas de movilidad sostenible que buscan reducir la congestión y las emisiones sin prohibir el acceso vehicular.
